\section{Related Works}
\label{sec:related_works}

\textbf{CAD datasets and representations.}
CAD datasets expose different aspects of parametric models: final geometry, construction history, or sketch-level constraints. ABC~\cite{Koch_2019_CVPR} provides large-scale 3D B-rep geometry, while DeepCAD~\cite{Wu_2021_ICCV}, Fusion 360 Gallery~\cite{willis2020fusion}, and recent extensions~\cite{rukhovich2025cad,fan2025history,lin2026unisketch,elistratov2026cadevolve} provide procedural CAD histories with increasing scale and command coverage. However, these resources provide limited supervision for design intent.

2D sketches expose design intent more directly because their constraints specify redesign behavior. SketchGraphs~\cite{seff2020sketchgraphs} provides large-scale Onshape sketches~\cite{onshape2026}, but most are under-constrained~\cite{casey2025aligning}; HistCAD~\cite{dong2025histcad} introduces a holistic constraint-aware representation, dataset, and editability benchmark for history-based CAD. Our dataset is designed specifically to capture valid design-intent diversity for the same geometry: we collect fully constrained sketches with three independent human annotations per sketch, enabling multi-reference evaluation rather than single-reference matching.

\textbf{CAD and sketch constraint generation.}
Most generative CAD methods produce command sequences or CAD programs from latent codes, text, images, point clouds, or engineering drawings~\cite{Wu_2021_ICCV,xu2022skexgen,xu2023hierarchical,khan2024cad,chen2025cadcrafter,rukhovich2025cad,kolodiazhnyi2025cadrille,li2025seek,elistratov2026cadevolve}. Their evaluations commonly emphasize execution validity, reconstruction fidelity, or static shape similarity. These criteria are important for geometric correctness, but they do not directly test whether the generated construction remains reusable under later parameter changes.

Sketch constraint generation addresses this issue more directly. Vitruvion~\cite{seff2021vitruvion} autoregressively generates sketch primitives and constraints, and later work extends this direction with richer representations, larger datasets, or multimodal conditioning~\cite{willis2021engineering,karadeniz2024davinci,wu2024cadvlm,mallis2025cad,lin2026unisketch,karadenizMiCADangelo}. However, standard evaluations still compare predictions against recorded SketchGraphs tokens, conflating command agreement with design-intent agreement. Casey~\etal~\cite{casey2025aligning} uses CAD solver feedback for reinforcement learning to improve sketch constrainedness, but solver rewards alone can overuse dimensions and reduce human preference. Thus, full constrainedness is necessary but not sufficient: sketch constraints should be evaluated for human-plausible design intent, not only solvability or single-reference token alignment.

\textbf{Design intent and editability.}
Design intent has long been studied in CAD as the modeling rationale that makes parametric models understandable, reusable, and robust under modification~\cite{iyer2006design,otey2018revisiting}. Importantly, prior work emphasizes that design intent is multifaceted: it may concern function, requirements, communication, redesign, and manufacturability, and may be expressed through sketch constraints, feature relationships, or modeling operations~\cite{otey2018revisiting}. This breadth has motivated studies on functional dimensioning, constraint taxonomies, robust modeling strategies, and constraint-based recognition of intended sketch relations~\cite{otto2015framework,naya2020role,nerenst2023parametric,kyratzi2021constraint}.

Due to the multifaceted nature, fully quantifying it with a single metric is not well-posed, so recent work operationalizes measurable facets. HistCAD~\cite{dong2025histcad}, for example, evaluates sketch editability by testing whether prescribed dimensional edits succeed while preserving specified constraints. We follow this behavior-based view for 2D sketch constraints, but characterize the full local space of admissible redesign motions rather than selected edit queries. This enables graded, syntax-invariant evaluation under multiple valid human intents.
